\section{Introduction}
A signalized intersection must divide service between competing traffic movements. Preset timing is straightforward to operate, but the allotted green time can become poorly matched to queues when directional demand changes. The timing manual describes the role of phase timing and traffic conditions in choosing a control strategy \cite{koonce}. An adaptive signal trial in Hong Kong and a separate city-scale study provide motivation for considering demand-responsive timing \cite{hk,wu}.

For a classroom prototype, the user needs to understand both why a plan was chosen and whether the controller applied it. The Traffic Timing Assistant addresses these needs at one intersection with north--south (NS) and east--west (EW) straight-through movements. Turning phases, pedestrian control and coordination across intersections are outside its scope.

The study asks whether demand-aware selection improves queue and waiting measures relative to fixed timing, whether a recommendation can be explained by its feature values and decision path, and whether the selected plan can be checked against device feedback. The contributions are a constrained tree-based selector over three plans, an ARM--FPGA division of decision and execution responsibilities, and an integrated workflow design for reporting data quality, decisions and readback together. The proposal reports existing model and control cores. The unified task interface, cross-source validation and common report are planned integration work.
